\documentclass[a4paper,12pt]{extarticle}
\usepackage{geometry}
\usepackage[T1]{fontenc}
\usepackage[utf8]{inputenc}
\usepackage[english,russian]{babel}
\usepackage{amsmath}
\usepackage{amsthm}
\usepackage{amssymb}
\usepackage{fancyhdr}
\usepackage{setspace}
\usepackage{graphicx}
\usepackage{colortbl}
\usepackage{tikz}
\usepackage{pgf}
\usepackage{subcaption}
\usepackage{listings}
\usepackage{indentfirst}
\usepackage[
    backend=biber,
    style=numeric,
    maxbibnames=99
]{biblatex}
\addbibresource{refs.bib}
\usepackage[colorlinks,citecolor=blue,linkcolor=blue,bookmarks=false,hypertexnames=true, urlcolor=blue]{hyperref}
\usepackage{indentfirst}
\usepackage{mathtools}
\usepackage{booktabs}
\usepackage[flushleft]{threeparttable}
\usepackage{tablefootnote}
\usepackage{longtable}

\usepackage{chngcntr} % нумерация графиков и таблиц по секциям
\counterwithin{table}{section}


\counterwithin{figure}{section}

\graphicspath{{graphics/}}%путь к рисункам

\makeatletter
% \renewcommand{\@biblabel}[1]{#1.} % Заменяем библиографию с квадратных скобок на точку:
\makeatother

\geometry{left=2.5cm}% левое поле
\geometry{right=1.0cm}% правое поле
\geometry{top=2.0cm}% верхнее поле
\geometry{bottom=2.0cm}% нижнее поле
\setlength{\parindent}{1.25cm}
\renewcommand{\baselinestretch}{1.5} % междустрочный интервал


\newcommand{\bibref}[3]{\hyperlink{#1}{#2 (#3)}} % biblabel, authors, year
\addto\captionsrussian{\def\refname{Список литературы (или источников)}}

\renewcommand{\theenumi}{\arabic{enumi}}% Меняем везде перечисления на цифра.цифра
\renewcommand{\labelenumi}{\arabic{enumi}}% Меняем везде перечисления на цифра.цифра
\renewcommand{\theenumii}{.\arabic{enumii}}% Меняем везде перечисления на цифра.цифра
\renewcommand{\labelenumii}{\arabic{enumi}.\arabic{enumii}.}% Меняем везде перечисления на цифра.цифра
\renewcommand{\theenumiii}{.\arabic{enumiii}}% Меняем везде перечисления на цифра.цифра
\renewcommand{\labelenumiii}{\arabic{enumi}.\arabic{enumii}.\arabic{enumiii}.}% Меняем везде перечисления на цифра.цифра

\begin{document}
\setcounter{page}{2}

{
	\hypersetup{linkcolor=black}
	\tableofcontents
}

\newpage

\newpage
\section*{Аннотация}

В рамках преддипломной практики выполнены работы по подготовке программного проекта веб-приложения для группового взаимодействия пользователей «Chattery».
Приложение предназначено для организации сообществ, обмена текстовыми сообщениями в тематических каналах, а также проведения групповых аудио- и видеозвонков.
Актуальность проекта обусловлена потребностью в доступной альтернативе существующим коммуникационным платформам. В ходе практики проведён анализ существующих решений, выявлены их функциональные ограничения, сформированы функциональные и нефункциональные требования к системе, а также описаны ключевые проектные решения для реализации серверной части и взаимодействия пользователей в реальном времени.
Результатом практики является уточнённая постановка задачи, обоснование необходимости разработки собственного решения, набор требований с прослеживаемостью от анализа аналогов и описание полученных проектных результатов.

\addcontentsline{toc}{section}{Аннотация}

\section*{Ключевые слова}
Веб-приложение, обмен сообщениями в реальном времени, видеосвязь, Go, JavaScript, PostgreSQL, Redis, WebSocket, WebRTC.
\pagebreak

\section{Введение}

Современные цифровые сообщества используют онлайн-платформы коммуникации для совместной работы, обсуждения общих интересов, проведения голосовых встреч и обмена сообщениями. Для подобных платформ важны не только отдельные функции чата или видеосвязи, но и возможность организовывать пользователей в устойчивые сообщества, разделять коммуникацию по тематическим каналам и поддерживать взаимодействие в режиме реального времени.

В рамках преддипломной практики рассматривается программный проект веб-прило\-жения «Chattery». Данная работа направлена на подготовку проектной основы: уточнение задачи, анализ аналогов, формирование требований, необходимых для разработки приложения.

\subsection{Актуальность темы}

Актуальность темы связана с необходимостью создания доступного веб-приложения для группового общения, которое может использоваться без зависимости от зарубежных платформ с ограниченной доступностью и без избыточной сложности интерфейса. Пользователю требуется единая среда, где можно создать сообщество, разделить обсуждения по темам, обмениваться сообщениями в реальном времени и подключаться к голосовым или видеозвонкам.

Существующие решения частично закрывают эту потребность, однако каждое из них имеет ограничения. Discord близок по функциональности, но не является гарантированно доступным решением для российского пользователя. Zoom и Jitsi Meet ориентированы прежде всего на разовые встречи, а не на долгосрочную структуру сообществ. TeamSpeak хорошо решает задачу голосовой связи, но не соответствует современным ожиданиям к веб-интерфейсу и видеосвязи. Matrix / Element предоставляет высокий уровень автономности и безопасности, однако требует более сложной настройки и менее интуитивен для массового пользователя.

\subsection{Постановка задачи}

Необходимо разработать веб-приложение для группового общения, в котором пользователи могут регистрироваться, создавать сообщества, организовывать тематические каналы, обмениваться сообщениями в реальном времени и участвовать в групповых аудио- и видеозвонках. В рамках практики основное внимание уделяется подготовке и обоснованию программного решения: анализу альтернатив, формированию требований, выбору технологий.

\subsection{Цель и задачи практики}

\textbf{Цель} -- подготовить и частично реализовать проектную основу веб-приложения «Chattery» для группового общения, включая анализ существующих решений, обоснование необходимости разработки собственного приложения и формирование требований.

\textbf{Задачи}:
\begin{enumerate}
    \item провести анализ существующих приложений для группового общения, видеоконференций и голосовой связи;
    \item выявить ограничения существующих решений, препятствующие их использованию как полноценной основы для целевого приложения;
    \item сформировать функциональные требования к веб-приложению на основе результатов анализа аналогов;
    \item сформировать нефункциональные требования к системе, включая требования к доступности, производительности, безопасности и удобству интерфейса;
    \item обосновать выбор основных технологий и методов реализации серверной части, обмена сообщениями и аудио-/видеосвязи;
\end{enumerate}

\newpage

\section{Анализ существующих приложений}

Ниже рассмотрены наиболее близкие по назначению платформы, применяемые для группового общения, организации сообществ и проведения голосовых или видеосеансов. Цель анализа состоит не только в перечислении возможностей аналогов, но и в выявлении причин, по которым требуется разработка собственного программного решения.

Для сравнения выбраны следующие критерии:
\begin{enumerate}
    \item наличие постоянных сообществ или серверов;
    \item поддержка тематических текстовых каналов;
    \item обмен сообщениями в реальном времени и хранение истории;
    \item поддержка групповой голосовой связи;
    \item поддержка групповой видеосвязи и демонстрации экрана;
    \item простота входа для пользователя;
    \item возможность самостоятельного развёртывания или контроля данных;
    \item доступность для целевой аудитории.
\end{enumerate}

\subsection{Discord}

Discord представляет собой универсальную платформу для группового общения, сочетающую текстовые каналы, голосовые и видеозвонки, а также развитые инструменты управления сообществами. Одной из ключевых особенностей сервиса является возможность создавать собственные серверы, внутри которых организуются тематические каналы для разных типов коммуникации. Платформа также поддерживает интеграцию ботов и сторонних сервисов, что существенно расширяет функциональность \cite{discord}.

С точки зрения функциональности Discord является наиболее близким аналогом разрабатываемого приложения: в нём есть серверы, каналы, история сообщений, голосовые комнаты и видеосвязь. Однако для целевого сценария он имеет два существенных ограничения. Во-первых, интерфейс и система настроек могут быть избыточными для пользователя, которому требуется простая коммуникационная среда. Во-вторых, доступность сервиса для российских пользователей ограничена \cite{discord_russia}. Следовательно, Discord показывает целевую модель функций, но не может быть принят как готовое решение для рассматриваемой задачи.

\subsection{Zoom}

Zoom является коммерческой платформой для видеоконференций, онлайн-встреч и вебинаров. Сервис рассчитан на проведение встреч с большим числом участников, поддерживает демонстрацию экрана, управление участниками, запись сеансов и другие инструменты, полезные в корпоративной среде \cite{zoom}.

Главное преимущество Zoom заключается в высокой стабильности видеосвязи и удобстве разовых мероприятий. Однако Zoom не является средой для постоянного формирования сообществ: в нём отсутствует привычная для Discord структура серверов и тематических каналов, а коммуникация строится вокруг конкретной встречи. Кроме того, для российских пользователей сохраняются риски ограничения доступа к сервису \cite{zoom_russia}. Таким образом, Zoom хорошо решает задачу видеоконференций, но не подходит как основа для приложения, где центральным объектом является сообщество с постоянными каналами и историей общения.

\subsection{Jitsi Meet}

Jitsi Meet -- это открытое решение для видеоконференций, которое позволяет быстро создавать групповые аудио- и видеозвонки прямо из браузера, без обязательной регистрации и установки дополнительного клиента. Платформа поддерживает шифрование, демонстрацию экрана и групповой чат, а также ориентирована на быстрый старт и простоту подключения \cite{jitsi}.

Сильной стороной Jitsi является минимальный барьер входа и возможность самостоятельного развёртывания. Однако сервис ориентирован прежде всего на видеовстречи, а не на длительное существование пользовательских сообществ. В нём отсутствует развитая модель серверов, тематических каналов, постоянной истории обсуждений и профилей участников. Следовательно, Jitsi может служить ориентиром для реализации браузерной видеосвязи, но не решает задачу создания полноценной коммуникационной платформы \cite{jitsi_limits}.

\subsection{Matrix / Element}

Matrix представляет собой децентрализованный протокол обмена сообщениями, основанный на федерации серверов и сквозном шифровании. На его основе построены клиенты, наиболее известным из которых является Element. Платформа поддерживает текстовые сообщения, передачу файлов, VoIP и видеосвязь \cite{matrix,element}.

Преимущество Matrix / Element состоит в автономности, контроле данных и развитой модели безопасности \cite{element_privacy}. Однако для массового пользователя такая система сложнее в освоении и настройке. В контексте данного проекта Matrix / Element является важным примером независимой коммуникационной инфраструктуры, но его сложность подтверждает необходимость более простого решения, ориентированного на базовые сценарии группового общения.

\subsection{TeamSpeak}

TeamSpeak -- это VoIP-клиент, применяемый прежде всего для голосового общения в игровых сообществах. Сервис работает на собственных или арендованных серверах и обеспечивает качественную голосовую связь, низкие задержки и гибкое управление правами доступа \cite{teamspeak}.

Ключевым ограничением TeamSpeak является узкая направленность на голосовое взаимодействие. В современных сценариях пользователям также требуются веб-интерфейс, текстовые каналы, видеосвязь, демонстрация экрана и удобная история сообщений. Поэтому TeamSpeak показывает важность голосовых комнат, но не подходит как универсальная платформа для разрабатываемого приложения.

\subsection*{Вывод по анализу альтернатив}
\addcontentsline{toc}{subsection}{Вывод по анализу альтернатив}

Проведённый анализ показывает, что ни одно из рассмотренных решений полностью не соответствует целевой постановке задачи. Вывод сравнительного анализа зафиксирован в таблице \ref{tab:analysis}.

Следовательно, разработка собственного веб-приложения обоснована необходимостью объединить в одном решении следующие свойства: доступность для целевой аудитории, структуру сообществ и тематических каналов, обмен сообщениями в реальном времени, групповые голосовые и видеозвонки, а также более простой пользовательский сценарий по сравнению с технически сложными или перегруженными аналогами.

\begin{table}
\centering
\small
\renewcommand{\arraystretch}{1.2}
\footnotesize
\begin{tabular}{|p{1.7cm}|p{1.3cm}|p{1.3cm}|p{1.3cm}|p{1.3cm}|p{1.3cm}|p{2.9cm}|p{2.9cm}|}
\hline
\textbf{Решение} & \textbf{Сооб\-щества} & \textbf{Текс\-товые каналы} & \textbf{Голос} & \textbf{Видео} & \textbf{Веб-доступ} & \textbf{Что отсутствует или ограничено} & \textbf{Вывод для проекта} \\
\hline
Discord & есть & есть & есть & есть & есть & ограниченная доступность в РФ, перегруженность интерфейса & нужен аналог с базовым набором функций и более простой моделью использования \\
\hline
Zoom & нет & нет & есть & есть & есть & нет постоянных сообществ, каналов и истории тематического общения & видеосвязь должна быть дополнением к сообществам, а не единственной функцией \\
\hline
Jitsi Meet & нет & нет & есть & есть & есть & нет развитых профилей, серверов, каналов и постоянной истории сообщений & полезна модель быстрого браузерного звонка, но нужны текстовые чаты \\
\hline
Matrix / Element & есть & есть & есть & есть & есть & высокая сложность настройки и освоения для массового пользователя & следует сохранить идею контроля данных, но упростить пользовательский сценарий \\
\hline
TeamSpeak & частично & нет & есть & ограни\-чено & ограни\-чено & устаревший интерфейс, слабая поддержка видео и текстового взаимодействия & голосовые комнаты важны, но должны быть встроены в современное приложение \\
\hline
\end{tabular}
\caption{Сравнение существующих решений по функциям, значимым для проекта}
\label{tab:analysis}
\end{table}

\newpage

\section{Формирование требований к приложению}

Требования к приложению сформированы на основе анализа альтернатив (таблица \ref{tab:analysis}). В таблице \ref{tab:req_trace} указано, какие ограничения существующих решений привели к появлению соответствующих требований.

\begin{table}[htbp]
\centering
\small
\renewcommand{\arraystretch}{1.2}
\footnotesize
\begin{tabular}{|p{5cm}|p{5cm}|p{6cm}|}
\hline
\textbf{Выявленное ограничение} & \textbf{Где проявляется} & \textbf{Требование к Chattery} \\
\hline
Нет единой структуры постоянных сообществ & Zoom, Jitsi Meet & реализовать серверы/сообщества как основные пространства взаимодействия \\
\hline
Нет тематических каналов внутри сообщества & Zoom, Jitsi Meet, TeamSpeak & реализовать текстовые и голосовые каналы внутри каждого сообщества \\
\hline
Сервис ориентирован только на видеовстречи или голосовую связь & Zoom, Jitsi Meet, TeamSpeak & совместить текстовое общение, голосовые и видеозвонки в одном веб-приложении \\
\hline
Недостаточная доступность для целевой аудитории & Discord, потенциально Zoom & обеспечить независимую реализацию, не зависящую от конкретной зарубежной платформы \\
\hline
Сложность настройки и освоения & Discord, Matrix / Element, TeamSpeak & сделать пользовательский интерфейс простым и сократить число действий для основных сценариев \\
\hline
Необходимость обмена событиями в реальном времени & все решения, поддерживающие чаты и звонки & использовать WebSocket для сообщений и событий звонков \\
\hline
Необходимость браузерной аудио- и видеосвязи & Discord, Zoom, Jitsi Meet, Matrix / Element & использовать WebRTC и реализовать серверный signaling \\
\hline
Необходимость сохранения истории общения & Discord, Matrix / Element & хранить сообщения в базе данных и поддерживать постраничную загрузку истории \\
\hline
\end{tabular}
\caption{Прослеживаемость требований от анализа аналогов}
\label{tab:req_trace}
\end{table}

\subsection{Функциональные требования}

К функциональным требованиям относятся требования, определяющие поведение системы и набор пользовательских возможностей.

\begin{enumerate}
    \item \textbf{Регистрация и аутентификация пользователей.} Система должна обеспечивать создание учётной записи, вход в систему и выход из неё.
    \item \textbf{Управление пользовательским профилем.} Пользователь должен иметь возможность просматривать и изменять профиль, включая имя, аватар и данные для аутентификации.
    \item \textbf{Создание сообществ.} Приложение должно поддерживать создание публичных и приватных сообществ, изменение их параметров, вступление пользователя в сообщество и выход из него. Требование сформировано из-за отсутствия полноценной модели постоянных сообществ в Zoom и Jitsi Meet.
    \item \textbf{Создание тематических каналов.} Внутри каждого сообщества должны поддерживаться тематические текстовые и голосовые каналы. Требование основано на сильной стороне Discord и ограничениях Zoom, Jitsi Meet и TeamSpeak, где такая структура отсутствует или выражена неполно.
    \item \textbf{Личные сообщения.} Система должна поддерживать личную переписку между двумя пользователями, список диалогов, отправку сообщений и хранение истории переписки.
    \item \textbf{Сообщения в каналах в реальном времени.} Передача текстовых сообщений должна осуществляться с использованием WebSocket. После отправки сообщение должно отображаться у других участников канала и сохраняться в базе данных.
    \item \textbf{История сообщений и пагинация.} Чат должен поддерживать загрузку истории сообщений постранично. Это необходимо для сохранения постоянного контекста общения, который отсутствует в сервисах, ориентированных только на разовые встречи.
    \item \textbf{Голосовые конференции.} Пользователи должны иметь возможность подключаться к групповым аудиозвонкам в голосовых каналах. Требование основано на сильных сторонах Discord и TeamSpeak.
    \item \textbf{Видеоконференции.} Приложение должно поддерживать групповые видеозвонки на базе WebRTC, включение и выключение камеры, а также возможность демонстрации экрана. Требование следует из возможностей Discord, Zoom и Jitsi Meet.
    \item \textbf{Сигналинг и управление звонками.} Для установления WebRTC-соединений необходимо реализовать обработку сигналов offer, answer и ICE candidates, а также события начала и завершения звонка, присоединения и выхода участников.
    \item \textbf{Отображение статуса участников.} Во время звонка необходимо отображать список активных участников и базовые состояния, например подключение, отключение микрофона или выход из звонка.
\end{enumerate}

\subsection{Нефункциональные требования}

Нефункциональные требования описывают ограничения и качественные характеристики системы.

\begin{enumerate}
    \item \textbf{Доступность.} Приложение должно быть реализовано как веб-система, доступная через браузер. Это снижает зависимость от установки отдельного клиента и учитывает ограничения, выявленные при анализе TeamSpeak.
    \item \textbf{Простота пользовательского интерфейса.} Основные действия пользователя -- регистрация, вход, переход к сообществу, отправка сообщения и подключение к звонку -- должны выполняться с минимальным числом шагов. Требование связано с избыточной сложностью Discord и Matrix / Element для части пользователей.
    \item \textbf{Производительность.} Система должна обеспечивать работу текстового чата и звонка для значительной группы пользователей, более десятка, без заметных задержек. Для этого целесообразно использовать WebSocket, Redis pub/sub и оптимизированные запросы к базе данных.
    \item \textbf{Масштабируемость.} Архитектура должна допускать дальнейшее увеличение числа пользователей, каналов и сообщений. Требование следует из ограничений решений, ориентированных только на разовые встречи, где постоянное накопление данных не является центральным сценарием.
    \item \textbf{Сохранность данных.} Сообщения, сведения о пользователях, сообществах и каналах должны храниться в PostgreSQL. Это необходимо для поддержки истории общения и восстановления состояния приложения после перезапуска.
    \item \textbf{Безопасность.} Пароли пользователей должны храниться в виде криптографических хэшей, а доступ к пользовательским действиям должен проверяться на серверной стороне.
    \item \textbf{Поддерживаемость.} Серверная часть должна быть разделена на логические компоненты, отвечающие за пользователей, сообщения, сообщества и звонки.
\end{enumerate}

\newpage

\section{Выбор методов и технологий решения}

Для реализации серверной части выбран язык Go, поскольку он подходит для разработки сетевых сервисов, хорошо поддерживает конкурентную обработку запросов и позволяет строить производительные backend-компоненты. Для хранения постоянных данных используется PostgreSQL: в ней целесообразно хранить пользователей, сообщества, каналы, сообщения и связи между сущностями.

Для обмена событиями в реальном времени выбран WebSocket. Он используется для доставки новых сообщений, обновления состояния каналов и передачи событий, связанных с подключением участников к звонкам. Redis применяется как вспомогательный компонент для кэширования сессий и реализации механизма публикации/подписки между backend-компонентами.

Для аудио- и видеосвязи выбран WebRTC, так как он позволяет устанавливать производительные мультимедийные соединения между браузерами пользователей и backend-сервисами. При этом серверная часть должна обеспечивать signaling: передачу offer, answer и ICE candidates между участниками звонка. Такой подход позволяет совместить браузерную доступность, работу в реальном времени и возможность дальнейшего масштабирования.

\section{Полученные результаты}

В ходе преддипломной практики были получены следующие результаты:
\begin{enumerate}
    \item уточнена актуальность и постановка задачи разработки веб-приложения «Chattery»;
    \item проведён анализ существующих решений: Discord,  Zoom, Jitsi Meet, Matrix / Element и TeamSpeak;
    \item выявлены ограничения аналогов, из-за которых требуется разработка собственного решения;
    \item подготовлена сводная таблица сравнения альтернатив;
    \item сформированы функциональные и нефункциональные требования к приложению;
    \item установлена прослеживаемость между недостатками аналогов и требованиями к системе;
    \item обоснован выбор технологий Go, PostgreSQL, Redis, WebSocket и WebRTC;
\end{enumerate}

\section{Заключение}

В ходе преддипломной практики выполнены работы, необходимые для подготовки программного проекта веб-приложения «Chattery». Были уточнены актуальность темы и постановка задачи, проведён анализ существующих коммуникационных платформ, выявлены их ограничения и сформированы требования к разрабатываемому приложению.

Основной вывод по результатам анализа состоит в том, что существующие решения закрывают отдельные части задачи, но не устраняют её полностью. Discord наиболее близок к целевой модели, однако имеет ограничения доступности и сложности интерфейса. Zoom и Jitsi Meet удобны для видеоконференций, но не предназначены для постоянных сообществ. TeamSpeak эффективен для голосовой связи, но не является современной универсальной веб-платформой. Matrix / Element обеспечивает автономность и безопасность, но сложнее для массового пользователя.

На основе выявленных ограничений обоснована необходимость разработки собственного приложения, совмещающего структуру сообществ, тематические каналы, сообщения в реальном времени, аудио- и видеосвязь, простоту интерфейса и возможность дальнейшего развития.

\newpage
\printbibliography[heading=bibintoc]

\newpage
\appendix

\end{document}
